% method_sim.tex
% Method: variance-corrected SIM composition.
% Six paragraphs ending with the Algorithm float:
%   (1) setup, SIM composition equation, and notation
%   (2) Gaussian-residual composition and why it fails
%   (3) generator-residual (naive) composition and its variance inflation
%   (4) variance correction + degenerate case + clipping rule
%   (5) R^2-preserving branch + SPY identity corollary
%   (6) algorithm intro + verification + contemporaneous-only caveat
%       (closes with the Algorithm float)

The HMM-WJ model produces a per-ticker marginal, but does not capture cross-asset dependence that is required for multi-asset applications. To generalize synthetic path generation beyond a single asset, we composed the per-ticker marginals using the single-index model (SIM) of Sharpe~\cite{sharpe1963}, which decomposes each asset's excess growth rate as:
\begin{equation}
  g_i(t) \;=\; \alpha_i \;+\; \beta_i\, g_m(t) \;+\; \eps_i(t),
  \qquad t = 1,\dots,T,
  \label{eq:sim-composition}
\end{equation}
where $g_m \in \R^T$ is the market growth-rate path with sample
variance $\sigma_m^2 = \Var(g_m)$, the pair $(\alpha_i, \beta_i)$
is the SIM intercept and slope for asset $i$, and $\eps_i$ is the
per-asset residual that the construction must specify (unrelated to
the jump probability $\epsilon$ of the marginal generator). Two
constraints govern the choice of $\eps_i$: the composed $g_i$ must
recover the calibrated $(\alpha_i, \beta_i)$ under ordinary least
squares (OLS), and it must inherit the per-asset marginal the HMM-WJ
generator was validated to reproduce (heavy tails, regime structure,
jumps).

A first option, \emph{Gaussian-residual composition},
fits a Normal distribution to the OLS residuals of the asset's real
history on $g_m$ and samples $\eps_i \sim N(0, \sigma_{\eps,i}^2)$ with
$\sigma_{\eps,i}^2 = (1 - R^2_{i,\real})\,\sigma_i^2$, where
$R^2_{i,\real}$ is the OLS coefficient of determination from
regressing the asset's history on $g_m$ and $\sigma_i^2$ is the
empirical variance of the asset's excess growth rates. This approach recovers
$(\alpha_i, \beta_i)$ and pins $\Var(g_i)$ to the empirical
$\sigma_i^2$, but the composed marginal becomes a mixture of a
Gaussian residual and a scaled market path. The asset's heavy tails,
regime structure, and jumps, which the per-asset HMM-WJ generator
was specifically built to reproduce, are absent from the resulting
$g_i$. Any fixed parametric residual family (Laplace, Student-$t$)
or empirical-bootstrap residual draw inherits the same limitation:
pinning the residual to a separate distribution discards the per-asset marginal we already have.

A second option uses the single-asset generator path $\tg_i$ (with sample variance $\sigma_{\gen,i}^2$) directly as the residual:
\begin{equation}\label{eq:naive-composition}
g_i = \alpha_i + \beta_i\, g_m + \tg_i
\end{equation}
This approach, which we call \emph{generator-residual composition}, captures the empirical per-asset marginal with variance
$\sigma_{\gen,i}^2 \approx \sigma_i^2$ calibrated to the asset, so its tails, regimes, and jumps are correct by construction. However, this approach also inflates the marginal variance. To see this, take the variance of both sides of Eq.~\eqref{eq:naive-composition}:
\begin{align}
  \Var(g_i)
   &= \Var(\alpha_i + \beta_i\, g_m + \tg_i)
       && (\text{naive composition}) \notag \\
   &= \Var(\beta_i\, g_m + \tg_i)
       && (\alpha_i \text{ is a constant}) \notag \\
   &= \beta_i^2\, \Var(g_m)
        + 2\beta_i\, \Cov(g_m,\, \tg_i)
        + \Var(\tg_i)
       && (\text{variance of a sum}) \notag \\
   &= \underbrace{\beta_i^2\, \sigma_m^2}_{\text{market}} \;+\; \underbrace{\sigma_{\gen,i}^2}_{\text{idiosyncratic}},
       && (\Cov(\tg_i, g_m) \approx 0)
   \label{eq:naive-inflation}
\end{align}
Thus, the variance of $g_i$ exceeds the per-asset (idiosyncratic) target $\Var(g_i) = \sigma_{\gen,i}^2$
by exactly $\beta_i^2 \sigma_m^2$, an inflation that grows
quadratically in the calibrated beta. 

To remove the inflation, we deflate the generator draw before composing:
form $\eps_i := s_i\,\tg_i$ for some scalar $s_i \in (0, 1]$ and
require $\Var(g_i) = \sigma_{\gen,i}^2$. Under the same independence
assumption $\Cov(\tg_i,\, g_m) \approx 0$,
$\Var(\beta_i\, g_m + s_i\,\tg_i) = \beta_i^2\, \sigma_m^2 +
s_i^2\, \sigma_{\gen,i}^2$; setting this equal to
$\sigma_{\gen,i}^2$ and solving yields the variance scale:
\begin{equation}
  \boxed{\;s_i^2 \;=\; 1 \;-\; \frac{\beta_i^2\, \sigma_m^2}{\sigma_{\gen,i}^2}\;}
  \label{eq:scale}
\end{equation}
Define the \emph{market loading ratio}:
\begin{equation}
  \rho_i \;\equiv\; \frac{\beta_i^2\, \sigma_m^2}{\sigma_{\gen,i}^2},
  \label{eq:rho}
\end{equation}
the share of the generator variance consumed by the market component;
then $s_i^2 = 1 - \rho_i$, valid whenever $\rho_i < 1$, with the
naive composition recovered in the $s_i = 1$ limit. Equation
\eqref{eq:scale} breaks down when $\rho_i \ge 1$: the market
loading alone would account for more than the entire generator
variance, leaving the idiosyncratic share non-positive, a case that arises
for high-$\beta_i$ assets paired with low-volatility generators or
whenever the synthetic market $g_m$ is more volatile than the market
on which $\beta_i$ was calibrated. To handle this case, we imposed a floor
$f \in (0, 1)$ on the idiosyncratic share by requiring $s_i^2 \ge f$,
and clipped $\beta_i$ downward whenever the floor was violated.
Writing the clip threshold as $\rho_i^{\max} = 1 - f$, the
piecewise clip rule is given by:
\begin{equation}
  \beta_i^{\eff} \;=\;
  \begin{cases}
    \beta_i & \text{if } \rho_i \le \rho_i^{\max} \\[4pt]
    \sign(\beta_i)\,\sqrt{(1 - f)\,\dfrac{\sigma_{\gen,i}^2}{\sigma_m^2}} & \text{if } \rho_i > \rho_i^{\max}
  \end{cases}
  \qquad
  s_i^2 \;=\;
  \begin{cases}
    1 - \rho_i & \text{if } \rho_i \le \rho_i^{\max} \\[4pt]
    f          & \text{if } \rho_i > \rho_i^{\max}
  \end{cases}
  \label{eq:clip}
\end{equation}
Here $\sign(\beta_i) \in \{-1, +1\}$ is the sign of the calibrated
$\beta_i$, so the clipped slope inherits the calibrated direction;
sign preservation keeps defensive (negative-$\beta$) assets
defensive even when clipped. The floor $f$ is a free parameter of
the construction; in our experiments we use $f = 0.10$, reserving at
least $10\%$ of each asset's variance for idiosyncratic noise.

Which target variance is appropriate depends on what the asset represents.
For a typical single stock the heavy tails, regime structure, and jumps live in the marginal that the HMM-WJ generator was built to reproduce. Thus, pinning $\Var(g_i)$ to $\sigma_{\gen,i}^2$ keeps the composition
aligned with the per-asset calibration. However, index-tracking ETFs reverse
the picture: what distinguishes a tracker is its OLS relationship to
the market, not its marginal. Against itself, SPY has $\beta = 1$,
$R^2 = 1$, and zero residual variance exactly; against SPY, QQQ has
$R^2 \approx 0.84$ and SPYG has $R^2 \approx 0.86$. A
variance-preserving composition applied to a tracker supplies a
residual whose variance bears no relation to the value
$R^2_{i,\real}$ implies, so OLS on the composed path returns
$\hat\beta$ and $\hat R^2$ that drift away from their calibrated
values. We therefore branch on the real-data calibration
$R^2_{i,\real}$: pick a threshold $R^2_{\mathrm{preserve}}$ that
separates index-tracking assets from single stocks; in our universe
the cumulative distribution sorts itself cleanly, with index ETFs
above $0.80$ and the most market-like single names below $0.65$
(Fig.~\ref{fig:r2_distribution}), so $R^2_{\mathrm{preserve}} = 0.80$
captures the trackers without sweeping in single names. For any asset with $R^2_{i,\real} \ge R^2_{\mathrm{preserve}}$,
inverting the population $R^2$ identity for
Eq.~\eqref{eq:sim-composition} for the residual variance that hits
the calibrated $R^2$ gives (Online
Appendix~\ref{sec:supp-sim-derivations}):
\begin{equation}
  \boxed{\;\sigma_{\eps,\mathrm{target}}^2 \;=\; \beta_i^2\, \sigma_m^2 \cdot
    \frac{1 - R^2_{i,\real}}{R^2_{i,\real}}\;}
  \label{eq:r2-target}
\end{equation}
The residual that achieves this target variance is the generator
path $\tg_i$ rescaled by the ratio of target to generator standard
deviation:
\begin{equation}
  \eps_i \;=\; \left(\sqrt{\frac{\sigma_{\eps,\mathrm{target}}^2}{\sigma_{\gen,i}^2}}\,\right)\tg_i,
  \label{eq:r2-rescale}
\end{equation}
which can then be copula-reordered and composed via
\eqref{eq:sim-composition}. The
construction recovers both calibrated quantities by design:
$\hat\beta \to \beta_i$ in the large-$T$ limit because $\eps_i$ is
uncorrelated with $g_m$ after the rank reorder, and $R^2 \to
R^2_{i,\real}$ by construction. Setting $R^2_{i,\real} = 1$ in
Eq.~\eqref{eq:r2-target} gives $\sigma_{\eps,\mathrm{target}}^2 = 0$,
so the SPY-against-itself identity $g_i = \alpha_i + \beta_i g_m$
falls out deterministically without special-casing; conversely, when
the rescale factor in Eq.~\eqref{eq:r2-rescale} exceeds one (which
happens when the synthetic market $g_m$ has lower variance than the
calibration market) the generator's tails are stretched to fit a
target larger than the original draw, and we detect and flag the case
so the loss of marginal fidelity is auditable.

The pieces combine into a single procedure (Algorithm~\ref{alg:hybrid}).
For each asset the calibrated $R^2_{i,\real}$ selects a branch:
assets with $R^2_{i,\real} \ge R^2_{\mathrm{preserve}}$ go through
the $R^2$-preserving construction of Eq.~\eqref{eq:r2-target}, and
the rest go through the variance-preserving construction of
Eq.~\eqref{eq:scale} with the clip of Eq.~\eqref{eq:clip} engaged
when the market loading ratio $\rho_i$ would otherwise exceed
$1 - f$. Each asset is labelled with the branch
that produced it ($R^2$-preserving, variance-preserving, or clipped
variance-preserving), recorded alongside the output so the path
taken per asset can be recovered without external bookkeeping. The
optional rank-reorder step in Algorithm~\ref{alg:hybrid} accepts
any joint dependence structure (empirical copula, t-copula, factor
copula); reordering
preserves the marginal variance of $\eps_i$, so the variance algebra
of Eq.~\eqref{eq:scale} is unaffected and whatever cross-sectional
structure is applied to $\eps_i$ passes through to $g_i$ unchanged.
By construction the composed $g_i$ satisfies three preservation
criteria in the large-$T$ limit: \emph{SIM recovery}, where OLS of
$g_i$ on $g_m$ returns $\hat\alpha_i = \alpha_i$ and
$\hat\beta_i = \beta_i^{\eff}$ (equal to $\beta_i$ on the unclipped
branch, attenuated by Eq.~\eqref{eq:clip} otherwise);
\emph{marginal fidelity}, where $\Var(g_i)$ matches the
branch-selected target ($\sigma_{\gen,i}^2$ on the
variance-preserving branch, or the value implied by the calibrated
$R^2_{i,\real}$ on the $R^2$-preserving branch); and
\emph{cross-sectional dependence}, where any joint dependence
structure applied to $\{\eps_i\}_{i=1}^N$ passes through unchanged
to $\{g_i\}_{i=1}^N$ by the linearity of
Eq.~\eqref{eq:sim-composition} in $\eps_i$. The verification
algebra, including the tail behavior implied by the floor $f$, is
laid out in Online Appendix~\ref{sec:supp-sim-derivations}. The
construction injects \emph{contemporaneous} market dependence only;
lagged dependence (lead-lag effects, market-driven volatility
clustering) requires a richer factor structure beyond the scope
here.

\begin{algorithm}[t]
\caption{Hybrid SIM composition with variance correction.}
\label{alg:hybrid}
\begin{algorithmic}[1]
  \REQUIRE Market growth-rate path $g_m \in \R^T$ with sample variance
    $\sigma_m^2$; per-asset generator growth-rate paths $\tg_i \in
    \R^T$ for $i = 1,\dots,N$;
    calibrated SIM parameters $(\alpha_i, \beta_i, R^2_{i,\real})$;
    floor $f \in (0,1)$; threshold $R^2_{\mathrm{preserve}} \in (0,1)$.
  \ENSURE Composed paths $g_i \in \R^T$ and per-asset construction flags.
  \FOR{each asset $i = 1,\dots,N$}
    \STATE Compute generator variance $\sigma_{\gen,i}^2 \gets \Var(\tg_i)$.
    \IF{$R^2_{i,\real} \ge R^2_{\mathrm{preserve}}$}
      \STATE $\sigma_{\eps,\mathrm{target}}^2 \gets \beta_i^2 \sigma_m^2 (1 - R^2_{i,\real})/R^2_{i,\real}$
        \hfill \COMMENT{Eq.~\eqref{eq:r2-target}}
      \STATE $\eps_i \gets \left(\sqrt{\sigma_{\eps,\mathrm{target}}^2 / \sigma_{\gen,i}^2}\,\right)\tg_i$
      \STATE $\beta_i^{\eff} \gets \beta_i$;\quad branch $\gets$ $R^2$-preserving
    \ELSE
      \STATE $\rho_i \gets \beta_i^2 \sigma_m^2 / \sigma_{\gen,i}^2$
        \hfill \COMMENT{Eq.~\eqref{eq:rho}}
      \IF{$\rho_i \le 1 - f$}
        \STATE $s_i^2 \gets 1 - \rho_i$;\quad $\beta_i^{\eff} \gets \beta_i$;\quad branch $\gets$ variance-preserving
      \ELSE
        \STATE $s_i^2 \gets f$;\quad
          $\beta_i^{\eff} \gets \sign(\beta_i)\sqrt{(1-f)\,\sigma_{\gen,i}^2/\sigma_m^2}$;\quad
          branch $\gets$ clipped variance-preserving
      \ENDIF
      \STATE $\eps_i \gets s_i\,\tg_i$
    \ENDIF
    \STATE \textit{(Optional)} apply cross-sectional rank reorder to $\eps_i$.
    \STATE $g_i \gets \alpha_i + \beta_i^{\eff}\, g_m + \eps_i$
      \hfill \COMMENT{Eq.~\eqref{eq:sim-composition}}
  \ENDFOR
\end{algorithmic}
\end{algorithm}
